\documentclass[10pt,a4paper]{scrartcl}

% ==============================================================================
% PACKAGES & SETUP
% ==============================================================================
\usepackage[utf8]{inputenc}
\usepackage[T1]{fontenc}
\usepackage{lmodern}
\usepackage[margin=13mm,top=14mm,bottom=15mm]{geometry}
\usepackage[table]{xcolor}
\usepackage{amsmath,amssymb}
\usepackage{booktabs}
\usepackage{tabularx}
\usepackage{multicol}
\usepackage{enumitem}
\usepackage{tikz}
\usetikzlibrary{positioning,calc,shapes.geometric,patterns,backgrounds}
\usepackage[many]{tcolorbox}
\usepackage{fontawesome5}
\usepackage{scrlayer-scrpage}

% ==============================================================================
% COLOR PALETTE (Modern Deep Teal & Sage Accent)
% ==============================================================================
\definecolor{TealDark}{RGB}{18,78,67}
\definecolor{TealMedium}{RGB}{35,110,96}
\definecolor{SageLight}{RGB}{234,244,240}
\definecolor{CoralAccent}{RGB}{217,107,39}
\definecolor{SlateDark}{RGB}{40,55,70}
\definecolor{BoxBg}{RGB}{249,251,250}
\definecolor{BorderGray}{RGB}{210,225,220}

% ==============================================================================
% HEADER & FOOTER CONFIGURATION (KOMA-Script)
% ==============================================================================
\clearpairofpagestyles
\addtokomafont{pageheadfoot}{\sffamily\color{SlateDark}}
\ihead{\textbf{\textcolor{TealDark}{Meridian Academy}} \quad\vline\quad Department of Biological Sciences}
\ohead{AP Biology \quad$\bullet$\quad Unit 4: Genetics}
\ifoot{\footnotesize Mendelian and Molecular Inheritance Worksheet}
\ofoot{\footnotesize Page \thepage\ of 3}

% ==============================================================================
% TCOLORBOX STYLES
% ==============================================================================
\tcbset{
  basebox/.style={
    enhanced,
    boxrule=0.8pt,
    arc=3pt,
    fonttitle=\bfseries\sffamily,
    attach boxed title to top left={xshift=10pt, yshift=-9pt},
    boxed title style={
      boxrule=0pt,
      arc=2pt,
      outer arc=2pt
    }
  }
}

\newtcolorbox{headerbox}[1][]{
  enhanced,
  colback=SageLight,
  colframe=TealDark,
  boxrule=1.2pt,
  arc=4pt,
  left=9pt, right=9pt, top=5pt, bottom=5pt,
  #1
}

\newtcolorbox{sectionbox}[2][]{
  basebox,
  colback=white,
  colframe=TealDark,
  coltitle=white,
  colbacktitle=TealDark,
  title={#2},
  left=8pt, right=8pt, top=7pt, bottom=6pt,
  #1
}

\newtcolorbox{workspace}[1][Work Area / Explanation]{
  enhanced,
  colback=BoxBg,
  colframe=TealMedium!40!white,
  boxrule=0.6pt,
  arc=3pt,
  title={\small\faPen{} #1},
  coltitle=TealDark,
  colbacktitle=SageLight,
  attach boxed title to top left={xshift=8pt, yshift=-8pt},
  boxed title style={boxrule=0.4pt, colframe=TealMedium!40!white},
  left=7pt, right=7pt, top=7pt, bottom=5pt
}

% Custom list styling
\setlist[enumerate,1]{label=\bfseries\color{TealDark}\arabic*., leftmargin=*, itemsep=2pt, topsep=1pt}
\setlist[enumerate,2]{label=\bfseries\color{TealMedium}(\alph*), leftmargin=*, itemsep=1pt, topsep=1pt}

\begin{document}

% Pre-load fontawesome font definitions to prevent tcolorbox font-loading issues
\sbox0{\faPen\faDna\faCheck}

% ==============================================================================
% DOCUMENT HEADER / METADATA BLOCK
% ==============================================================================
\begin{headerbox}
\begin{tabularx}{\textwidth}{@{} X r @{}}
  {\Large\textbf{\textcolor{TealDark}{\sffamily UNIT 4: GENETICS \& PATTERNS OF INHERITANCE}}} & \textcolor{TealDark}{\Large\faDna} \\
  {\small\sffamily Problem Set, Punnett Square Analysis \& Pedigree Interpretation} & {\small\sffamily Academic Year 2025--2026}
\end{tabularx}
\vspace{2pt}
\textcolor{TealMedium}{\rule{\textwidth}{0.8pt}}
\vspace{2pt}
\sffamily
\begin{tabularx}{\textwidth}{@{} l X l X l X @{}}
  \textbf{Student Name:} & \underline{\hspace{3.2cm}} & \textbf{Date:} & \underline{\hspace{1.6cm}} & \textbf{Period:} & \underline{\hspace{0.8cm}} \\
  \textbf{Instructor:} & Dr. E. Vance & \textbf{Total Score:} & \underline{\quad / 50 pts} & \textbf{Grade:} & \underline{\hspace{0.8cm}}
\end{tabularx}
\end{headerbox}

\vspace{4pt}

% ==============================================================================
% SECTION 1: FOUNDATIONS & TERMINOLOGY
% ==============================================================================
\begin{sectionbox}{Section 1: Fundamental Principles \& Terminology (8 Points)}
\sffamily
Complete the table below by providing the correct genetic term, symbol notation, or description.

\vspace{1pt}
\rowcolors{2}{SageLight!50}{white}
\begin{tabularx}{\textwidth}{|>{\bfseries\color{TealDark}}p{3.0cm}|X|>{\centering\arraybackslash}p{2.4cm}|>{\centering\arraybackslash}p{2.3cm}|}
\hline
\rowcolor{TealDark}
\color{white}\textbf{Genetic Term} & \color{white}\textbf{Definition \& Biological Significance} & \color{white}\textbf{Example Alleles} & \color{white}\textbf{Phenotype Outcome} \\
\hline
Homozygous Dominant & Possesses two identical dominant alleles for a given gene locus. & $BB$ & Dominant Trait \\
\hline
Heterozygous & Possesses two distinct alleles for a single gene locus. & \underline{\hspace{1.5cm}} & Dominant Trait \\
\hline
Phenotype & \underline{\hspace{5.0cm}} & Physical Appearance & Observable Trait \\
\hline
Codominance & \underline{\hspace{5.5cm}} & $I^A I^B$ & \underline{\hspace{1.5cm}} \\
\hline
Recessive Allele & Allele whose phenotype is masked in the presence of a dominant allele. & $b$ & Expressed if $bb$ \\
\hline
\end{tabularx}
\end{sectionbox}

\vspace{4pt}

% ==============================================================================
% SECTION 2: MONOHYBRID CROSSES
% ==============================================================================
\begin{sectionbox}{Section 2: Monohybrid Inheritance \& Punnett Square Analysis (10 Points)}
\begin{enumerate}
  \item \textbf{Pod Shape Inheritance in Garden Peas (\textit{Pisum sativum}):}
  In pea plants, inflated pod shape ($I$) is completely dominant over constricted pod shape ($i$). A heterozygous inflated-pod plant is crossed with a homozygous constricted-pod plant.

  \vspace{1pt}
  \begin{minipage}[t]{0.56\textwidth}
    \begin{enumerate}
      \item State parental genotypes:
      \begin{itemize}[label={--}, leftmargin=10pt, itemsep=0pt]
        \item Parent 1 (Inflated): \underline{\hspace{1.5cm}}
        \item Parent 2 (Constricted): \underline{\hspace{1.4cm}}
      \end{itemize}
      \item Complete the Punnett Square on the right.
      \item Calculate phenotypic ratio of offspring:\\
      Inflated : Constricted = \underline{\hspace{1.5cm}}
      \item If 480 offspring seeds are collected, calculate expected count of constricted pod plants:\\
      Expected Count = \underline{\hspace{1.5cm}} plants.
    \end{enumerate}
  \end{minipage}%
  \hfill
  \begin{minipage}[t]{0.40\textwidth}
    \centering
    \begin{tcolorbox}[colback=white, colframe=TealMedium, boxrule=0.8pt, title={\centering\small\sffamily\bfseries Monohybrid Cross Punnett Square}, coltitle=white, width=\linewidth, top=1pt, bottom=1pt]
      \centering
      \vspace{1pt}
      \begin{tikzpicture}[scale=0.78]
        \draw[thick, TealDark] (0,0) grid (3,3);
        \draw[very thick, TealDark] (0,0) rectangle (3,3);
        \draw[line width=1.2pt, TealDark] (0,1.5) -- (3,1.5);
        \draw[line width=1.2pt, TealDark] (1.5,0) -- (1.5,3);
        
        % Maternal Gametes
        \node[font=\bfseries\sffamily\Large\color{CoralAccent}] at (0.75, 3.35) {$I$};
        \node[font=\bfseries\sffamily\Large\color{CoralAccent}] at (2.25, 3.35) {$i$};
        \node[font=\sffamily\footnotesize\bfseries\color{TealDark}] at (1.5, 3.75) {Maternal Gametes};
        
        % Paternal Gametes
        \node[font=\bfseries\sffamily\Large\color{TealDark}] at (-0.4, 2.25) {$i$};
        \node[font=\bfseries\sffamily\Large\color{TealDark}] at (-0.4, 0.75) {$i$};
        \node[font=\sffamily\footnotesize\bfseries\color{TealDark}, rotate=90] at (-0.8, 1.5) {Paternal Gametes};
        
        % Offspring genotypes
        \node[font=\bfseries\sffamily\large\color{SlateDark}] at (0.75, 2.25) {$Ii$};
        \node[font=\bfseries\sffamily\large\color{SlateDark}] at (2.25, 2.25) {$ii$};
        \node[font=\bfseries\sffamily\large\color{SlateDark}] at (0.75, 0.75) {$Ii$};
        \node[font=\bfseries\sffamily\large\color{SlateDark}] at (2.25, 0.75) {$ii$};
      \end{tikzpicture}
      \vspace{1pt}
    \end{tcolorbox}
  \end{minipage}

  \vspace{2pt}

  \item \textbf{Coat Color Genetics in Rex Rabbits:}
  In Rex rabbits, black coat color ($B$) is dominant to chinchilla gray ($b$). A geneticist crosses two black rabbits. The resulting litter consists of 27 black kits and 9 chinchilla gray kits.

  \begin{enumerate}
    \item Determine parental genotypes and justify your answer using Mendelian ratios.
    \item Perform the necessary Punnett square analysis to support your conclusion.
  \end{enumerate}

  \vspace{1pt}
  \begin{workspace}[Deduction \& Genetic Proof]
    \vspace{0.9cm}
  \end{workspace}
\end{enumerate}
\end{sectionbox}

\newpage

% ==============================================================================
% SECTION 3: DIHYBRID CROSSES & INDEPENDENT ASSORTMENT
% ==============================================================================
\begin{sectionbox}{Section 3: Dihybrid Inheritance \& Independent Assortment (12 Points)}
\begin{enumerate}
  \item \textbf{\textit{Drosophila melanogaster} Dihybrid Cross Analysis:}
  In fruit flies, gray body color ($E$) is dominant over ebony body ($e$), and normal wings ($V$) are dominant over vestigial wings ($v$). A male fly heterozygous for both traits ($EeVv$) is mated with a female fly also heterozygous for both traits ($EeVv$).

  \begin{enumerate}
    \item List the four possible gamete combinations produced by each parent fly:
    \begin{center}
      Gametes: \quad \fbox{\hspace{6pt}$EV$\hspace{6pt}} \quad \fbox{\hspace{6pt}$Ev$\hspace{6pt}} \quad \fbox{\hspace{6pt}$eV$\hspace{6pt}} \quad \fbox{\hspace{6pt}$ev$\hspace{6pt}}
    \end{center}
    \item Complete the 4x4 Dihybrid Punnett Square below by writing the resulting 4-allele genotype in each cell.
  \end{enumerate}

  \vspace{1pt}
  \begin{center}
  \begin{tcolorbox}[colback=white, colframe=TealDark, boxrule=1pt, title={\centering\sffamily\bfseries Dihybrid Cross Grid ($EeVv \times EeVv$)}, coltitle=white, width=0.82\textwidth, top=1pt, bottom=1pt]
    \centering
    \vspace{1pt}
    \begin{tikzpicture}[scale=0.76]
      % Grid boundaries
      \draw[thick, TealDark] (0,0) grid (6.4,6.4);
      \draw[very thick, TealDark] (0,0) rectangle (6.4,6.4);
      \foreach \x in {1.6, 3.2, 4.8} {
        \draw[line width=1.1pt, TealDark] (\x,0) -- (\x,6.4);
        \draw[line width=1.1pt, TealDark] (0,\x) -- (6.4,\x);
      }
      
      % Top Maternal Gametes
      \node[font=\bfseries\sffamily\color{CoralAccent}] at (0.8, 6.75) {$EV$};
      \node[font=\bfseries\sffamily\color{CoralAccent}] at (2.4, 6.75) {$Ev$};
      \node[font=\bfseries\sffamily\color{CoralAccent}] at (4.0, 6.75) {$eV$};
      \node[font=\bfseries\sffamily\color{CoralAccent}] at (5.6, 6.75) {$ev$};
      \node[font=\sffamily\small\bfseries\color{TealDark}] at (3.2, 7.2) {Female Gametes ($EeVv$)};
      
      % Left Paternal Gametes
      \node[font=\bfseries\sffamily\color{TealDark}] at (-0.5, 5.6) {$EV$};
      \node[font=\bfseries\sffamily\color{TealDark}] at (-0.5, 4.0) {$Ev$};
      \node[font=\bfseries\sffamily\color{TealDark}] at (-0.5, 2.4) {$eV$};
      \node[font=\bfseries\sffamily\color{TealDark}] at (-0.5, 0.8) {$ev$};
      \node[font=\sffamily\small\bfseries\color{TealDark}, rotate=90] at (-0.95, 3.2) {Male Gametes ($EeVv$)};

      % Pre-filled reference genotypes & blanks
      \node[font=\sffamily\small\bfseries\color{SlateDark}] at (0.8, 5.6) {$EEVV$};
      \node[font=\sffamily\small\bfseries\color{SlateDark}] at (2.4, 5.6) {$EEVv$};
      \node[font=\sffamily\small\bfseries\color{SlateDark}] at (4.0, 5.6) {$EeVV$};
      \node[font=\sffamily\small\bfseries\color{SlateDark}] at (5.6, 5.6) {$EeVv$};

      \node[font=\sffamily\small\bfseries\color{SlateDark}] at (0.8, 4.0) {$EEVv$};
      \node[font=\sffamily\small\bfseries\color{SlateDark}] at (2.4, 4.0) {$EEvv$};
      \node[font=\sffamily\small\bfseries\color{SlateDark}] at (4.0, 4.0) {$EeVv$};
      \node[font=\sffamily\small\color{gray!70}] at (5.6, 4.0) {$Eevv$};

      \node[font=\sffamily\small\bfseries\color{SlateDark}] at (0.8, 2.4) {$EeVV$};
      \node[font=\sffamily\small\bfseries\color{SlateDark}] at (2.4, 2.4) {$EeVv$};
      \node[font=\sffamily\small\color{gray!70}] at (4.0, 2.4) {$eeVV$};
      \node[font=\sffamily\small\color{gray!70}] at (5.6, 2.4) {$eeVv$};

      \node[font=\sffamily\small\bfseries\color{SlateDark}] at (0.8, 0.8) {$EeVv$};
      \node[font=\sffamily\small\color{gray!70}] at (2.4, 0.8) {$Eevv$};
      \node[font=\sffamily\small\color{gray!70}] at (4.0, 0.8) {$eeVv$};
      \node[font=\sffamily\small\bfseries\color{TealDark}] at (5.6, 0.8) {$eevv$};
    \end{tikzpicture}
    \vspace{1pt}
  \end{tcolorbox}
  \end{center}

  \vspace{1pt}
  \begin{enumerate}[resume]
    \item Summarize expected phenotypic counts out of 16 total offspring:
    \begin{multicols}{2}
      \begin{itemize}[label={\small\faCheck}]
        \item Gray Body, Normal Wings: \underline{\hspace{1.2cm}} / 16
        \item Gray Body, Vestigial Wings: \underline{\hspace{1.2cm}} / 16
        \item Ebony Body, Normal Wings: \underline{\hspace{1.2cm}} / 16
        \item Ebony Body, Vestigial Wings: \underline{\hspace{1.2cm}} / 16
      \end{itemize}
    \end{multicols}

    \item \textbf{Meiotic Mechanism:} Explain how the chromosome alignment during \textbf{Metaphase I} of meiosis produces independent assortment of unlinked genes.
  \end{enumerate}

  \vspace{1pt}
  \begin{workspace}[Biological Explanation]
    \vspace{0.8cm}
  \end{workspace}
\end{enumerate}
\end{sectionbox}

\vspace{3pt}

% ==============================================================================
% SECTION 4: NON-MENDELIAN INHERITANCE
% ==============================================================================
\begin{sectionbox}{Section 4: Complex \& Non-Mendelian Patterns (10 Points)}
\begin{enumerate}
  \item \textbf{Incomplete Dominance in Snapdragons (\textit{Antirrhinum majus}):}
  Red flower color is governed by allele $C^R$ and white flower color by allele $C^W$. Heterozygous plants ($C^R C^W$) exhibit pink flowers.
  \begin{enumerate}
    \item A pink snapdragon is crossed with a white snapdragon. Determine expected genotypic and phenotypic ratios of offspring.
  \end{enumerate}

  \vspace{1pt}
  \begin{workspace}[Snapdragon Cross Work Area]
    \vspace{0.8cm}
  \end{workspace}

  \vspace{2pt}
  \item \textbf{Multiple Alleles \& Codominance: Human ABO Blood Types:}
  Elena (Blood Type A) and Marcus (Blood Type B) have a child, Julian, who has Blood Type O ($ii$).
  \begin{enumerate}
    \item Identify the precise genotypes of Elena: \underline{\hspace{1.8cm}} and Marcus: \underline{\hspace{1.8cm}}.
    \item What is the probability that their next child will have Blood Type AB? Show your cross.
    \item \textbf{Immunohematology:} Why can an individual with Type O blood safely donate red blood cells to any recipient regardless of blood type?
  \end{enumerate}
\end{enumerate}
\end{sectionbox}

\newpage

% ==============================================================================
% SECTION 5: SEX-LINKED TRAITS & PEDIGREE ANALYSIS
% ==============================================================================
\begin{sectionbox}{Section 5: Sex-Linked Inheritance \& Pedigree Analysis (10 Points)}
\sffamily
The pedigree below traces the inheritance of \textbf{Hemophilia A} (an X-linked recessive bleeding disorder caused by a deficiency in Factor VIII) across three generations of the Nexa family.

\vspace{2pt}
\begin{center}
\begin{tikzpicture}[
  male/.style={rectangle, draw=TealDark, thick, minimum size=5.5mm, inner sep=0pt},
  female/.style={circle, draw=TealDark, thick, minimum size=5.5mm, inner sep=0pt},
  affected/.style={fill=TealDark!85, text=white},
  carrier/.style={fill=CoralAccent!45},
  unaffected/.style={fill=white},
  line/.style={draw=TealDark, thick},
  labfont/.style={font=\sffamily\small\bfseries\color{SlateDark}}
]
  % Generation Labels
  \node[labfont] at (-1.5, 0) {Gen I};
  \node[labfont] at (-1.5, -1.5) {Gen II};
  \node[labfont] at (-1.5, -3.0) {Gen III};

  % --- Generation I ---
  \node[male, unaffected, label=below:{\footnotesize I-1}] (I1) at (1.2, 0) {};
  \node[female, carrier, label=below:{\footnotesize I-2}] (I2) at (4.2, 0) {};
  \draw[line] (I1) -- (I2);
  \coordinate (M1) at (2.7, 0);

  % --- Generation II ---
  \coordinate (P2) at (2.7, -0.75);
  \draw[line] (M1) -- (P2);
  \draw[line] (0.2, -0.75) -- (5.2, -0.75);
  \draw[line] (0.2, -0.75) -- (0.2, -1.5);
  \draw[line] (1.87, -0.75) -- (1.87, -1.5);
  \draw[line] (3.53, -0.75) -- (3.53, -1.5);
  \draw[line] (5.2, -0.75) -- (5.2, -1.5);

  \node[female, unaffected, label=below:{\footnotesize II-1}] (II1) at (0.2, -1.5) {};
  \node[male, affected, label=below:{\footnotesize II-2}] (II2) at (1.87, -1.5) {};
  \node[female, carrier, label=below:{\footnotesize II-3}] (II3) at (3.53, -1.5) {};
  \node[male, unaffected, label=below:{\footnotesize II-4}] (II4) at (5.2, -1.5) {};

  % Marriage II-3 with II-5
  \node[male, unaffected, label=below:{\footnotesize II-5}] (II5) at (7.0, -1.5) {};
  \draw[line] (II3) -- (II5);
  \coordinate (M2) at (5.265, -1.5);

  % --- Generation III ---
  \coordinate (P3) at (5.265, -2.25);
  \draw[line] (M2) -- (P3);
  \draw[line] (3.8, -2.25) -- (6.7, -2.25);
  \draw[line] (3.8, -2.25) -- (3.8, -3.0);
  \draw[line] (5.25, -2.25) -- (5.25, -3.0);
  \draw[line] (6.7, -2.25) -- (6.7, -3.0);

  \node[male, affected, label=below:{\footnotesize III-1}] (III1) at (3.8, -3.0) {};
  \node[female, carrier, label=below:{\footnotesize III-2}] (III2) at (5.25, -3.0) {};
  \node[male, unaffected, label=below:{\footnotesize III-3}] (III3) at (6.7, -3.0) {};

  % --- Legend Box ---
  \begin{scope}[shift={(8.2, 0.2)}]
    \draw[thick, TealDark, fill=SageLight, rounded corners=4pt] (-0.3, 0.3) rectangle (4.2, -3.4);
    \node[anchor=west, font=\sffamily\bfseries\color{TealDark}] at (0, -0.1) {Pedigree Chart Legend};
    \node[male, unaffected, scale=0.6] at (0.3, -0.7) {};
    \node[anchor=west, font=\sffamily\small] at (0.8, -0.7) {Unaffected Male};
    \node[female, unaffected, scale=0.6] at (0.3, -1.3) {};
    \node[anchor=west, font=\sffamily\small] at (0.8, -1.3) {Unaffected Female};
    \node[male, affected, scale=0.6] at (0.3, -1.9) {};
    \node[anchor=west, font=\sffamily\small] at (0.8, -1.9) {Hemophilic Individual};
    \node[female, carrier, scale=0.6] at (0.3, -2.5) {};
    \node[anchor=west, font=\sffamily\small] at (0.8, -2.5) {Heterozygous Carrier};
    \node[anchor=west, font=\sffamily\scriptsize\color{TealMedium}] at (0.0, -3.1) {$X^H = \text{Normal}, X^h = \text{Hemophilia}$};
  \end{scope}
\end{tikzpicture}
\end{center}

\vspace{1pt}

\begin{enumerate}
  \item \textbf{Genotype Determination:} Assign the most precise genotype for each individual:
  \begin{multicols}{2}
    \begin{itemize}[label={\small\faDna}]
      \item Individual I-1: \underline{\hspace{2.2cm}}
      \item Individual I-2: \underline{\hspace{2.2cm}}
      \item Individual II-2: \underline{\hspace{2.2cm}}
      \item Individual II-3: \underline{\hspace{2.2cm}}
      \item Individual III-1: \underline{\hspace{2.2cm}}
      \item Individual III-2: \underline{\hspace{2.2cm}}
    \end{itemize}
  \end{multicols}

  \item \textbf{Probability Calculation:} If Individual III-2 (carrier female) marries an unaffected male ($X^H Y$), calculate the probability of the following outcomes:
  \begin{enumerate}
    \item Probability that their first \textbf{son} will have hemophilia: \underline{\hspace{2.0cm}}
    \item Probability that their first \textbf{daughter} will be a carrier: \underline{\hspace{2.0cm}}
    \item Overall probability of having an affected child of \textbf{any gender}: \underline{\hspace{2.0cm}}
  \end{enumerate}

  \item \textbf{Genetic Reasoning:} Why are X-linked recessive traits expressed significantly more frequently in males ($XY$) than in females ($XX$)?
\end{enumerate}

\vspace{1pt}
\begin{workspace}[X-Linked Inheritance Mechanism]
  \vspace{0.8cm}
\end{workspace}
\end{sectionbox}

\vspace{3pt}

% ==============================================================================
% SECTION 6: CHI-SQUARE ANALYSIS
% ==============================================================================
\begin{sectionbox}{Section 6: Statistical Analysis -- Chi-Square ($\chi^2$) Goodness-of-Fit (10 Points)}
\sffamily
A research team at \textbf{Vertex Genetics Institute} crossed two heterozygous tall, purple-flowered pea plants ($TtPp \times TtPp$). Out of a total sample of \textbf{800 offspring plants}, the observed phenotypic distribution was recorded in the table below.

Test the null hypothesis ($H_0$) that the data fits the expected Mendelian \textbf{9:3:3:1} dihybrid ratio at a significance level of $p = 0.05$.

\vspace{2pt}
\rowcolors{2}{SageLight!40}{white}
\begin{tabularx}{\textwidth}{|X|>{\centering\arraybackslash}p{2.0cm}|>{\centering\arraybackslash}p{2.0cm}|>{\centering\arraybackslash}p{2.0cm}|>{\centering\arraybackslash}p{2.6cm}|}
\hline
\rowcolor{TealDark}
\color{white}\textbf{Phenotype Category} & \color{white}\textbf{Observed ($O$)} & \color{white}\textbf{Expected Ratio} & \color{white}\textbf{Expected ($E$)} & \color{white}\textbf{$(O - E)^2 / E$} \\
\hline
Tall, Purple Flowers & 462 & $9/16$ & 450 & $(462 - 450)^2 / 450 = 0.320$ \\
\hline
Tall, White Flowers & 138 & $3/16$ & 150 & \underline{\hspace{1.6cm}} \\
\hline
Dwarf, Purple Flowers & 154 & $3/16$ & 150 & \underline{\hspace{1.6cm}} \\
\hline
Dwarf, White Flowers & 46 & $1/16$ & 50 & \underline{\hspace{1.6cm}} \\
\hline
\multicolumn{4}{|r|}{\textbf{Total Chi-Square Value ($\chi^2 = \sum \frac{(O-E)^2}{E}$):}} & \underline{\hspace{1.8cm}} \\
\hline
\end{tabularx}

\vspace{2pt}
\begin{enumerate}
  \item Degrees of freedom ($df = k - 1$): \underline{\hspace{1.6cm}} \quad (Critical value at $p=0.05$ with $df=3$ is \textbf{7.82}).
  \item \textbf{Conclusion Statement:} Do you accept or reject the Null Hypothesis ($H_0$)? State your biological rationale.
\end{enumerate}

\vspace{1pt}
\begin{workspace}[Chi-Square Decision \& Rationale]
  \vspace{0.8cm}
\end{workspace}
\end{sectionbox}

\end{document}
